\documentclass[11pt,a4paper]{article}
\usepackage[T1]{fontenc}
\usepackage[utf8]{inputenc}
\usepackage[english]{babel}
\usepackage{amsmath,amssymb}
\usepackage[margin=2cm]{geometry}
\setlength{\parindent}{0pt}
\setlength{\emergencystretch}{3em}
\allowdisplaybreaks

\title{About \texttt{book\_integrals.tex}}
\author{}
\date{}
\begin{document}
\maketitle
This note explains what the test data in \texttt{book\_integrals.tex} is used for. That file is input for the engine in \texttt{main.c}; it is read by the test \texttt{test\_integrals.c} in this folder.

\begin{enumerate}
\item \textbf{Purpose.} The file holds 62 true statements. Each says that an antiderivative printed in a textbook is right, $\int f \, dx = F + c$, or that a definite integral has the printed value, $\int_{a}^{b} f \, dx = V$. An antiderivative statement means: on every open interval where $F$ and $f$ are defined, $F' = f$.

\item \textbf{Where the data comes from.} The exercises are those of \texttt{photos/dati\_esercizi/integrali.jsonl} (university and school sheets: Firenze, Verona, Ferrara, Trieste, De Agostini Scuola, Brescia, Colorado, edutecnica), whose printed results an independent numerical oracle confirmed. In two exercises the printed result is wrong (firenze-5d lacks the coefficients $6$ and $3$, ferrara-R7 has $-\ln$ where $-4\ln$ is right): the statement here has the corrected result, and the printed one is in \texttt{false\_integrals.tex}.

\item \textbf{Notation changed from the sources.} The university sources write $\log$ for the natural logarithm; the reader of the program reads $\log$ without a base as the logarithm in base $10$, as school books do, so here every such $\log$ is written $\ln$. Ferrara writes $\operatorname{arctg}$, which the reader does not know; here it is $\arctan$. A bare bar, $|x - 2|$, is read as an absolute value (paragraph 9c of the reader).

\item \textbf{What it exercises.} The verification of an antiderivative (rules 30a and 30b): the derivative of $F$ by the rules of chapter 27, $F' - f$ reduced to zero by exact algebra (chapter 12), and the derivative read from the Taylor series of $F$ with interval arithmetic at sample points. For the definite integrals, the search for an antiderivative (30c--30h: partial fractions, polynomials times exponentials, integration by parts, substitution), the conditions of the fundamental theorem checked with interval arithmetic on a subdivision of $[a, b]$, the split at the zeros of an absolute value (30i-2), the exact values of sines, cosines and arctangents at multiples of $\frac{\pi}{12}$, and the Riemann enclosure of the integral as the independent check (30i).

\item \textbf{Variety.} Rational functions with distinct, repeated and complex roots and with repeated quadratic factors; logarithms with and without absolute values; integration by parts, cyclic ones included; substitutions with roots, exponentials and logarithms; powers and products of sines and cosines; Weierstrass substitutions; square roots of quadratics (checked, not found); definite integrals with $\pi$ and $e$ in the bounds and in the values, with $|x - 1|$, $|\sin x|$ and $|6x - \pi|$ in the integrand, and one of size $6.5 \cdot 10^{12}$.

\item \textbf{Expected result.} The line \texttt{PASS: 62 of 62 statements proved, each with a written proof}.

\item \textbf{How to run it.} From the top folder of the repository:
\begin{verbatim}
cc -std=c11 -O2 -Wall -Wextra -o /tmp/test_integrals \
    tests/integrals/test_integrals.c && /tmp/test_integrals
\end{verbatim}
\end{enumerate}
\end{document}
